% Rédaction : KOUMTOUBRE Davy ; relecture : BADO Lionel
\chapter{API REST de transactions}\label{chap:api}

\section{Responsabilités et organisation}
L'API expose les transactions conservées dans PostgreSQL et les résultats du chargement KNIME.
Elle est construite avec Express~5 et séparée en couches afin que les routes ne portent ni les
règles métier ni les détails SQL :
\begin{enumerate}
  \item les routes HTTP reçoivent la requête et choisissent le code de réponse ;
  \item les modules de validation normalisent et contrôlent les entrées ;
  \item le dépôt de transactions exécute les opérations sur PostgreSQL ;
  \item un gestionnaire central transforme les erreurs en réponses JSON cohérentes.
\end{enumerate}
Cette séparation permet de tester les routes avec un dépôt en mémoire, puis de vérifier les
requêtes SQL séparément contre PostgreSQL (chapitre~\ref{chap:tests}).

\section{Routes et opérations}
Les transactions utilisent leur référence métier \texttt{TXN-xxxxxx} dans l'URL. Cette clé reste
stable entre les chargements, contrairement à un identifiant technique de base.

\begin{table}[H]
\centering\small
\begin{tabularx}{\linewidth}{@{}l l X@{}}
\toprule
\textbf{Méthode} & \textbf{Route} & \textbf{Comportement principal} \\
\midrule
GET & \texttt{/api/transactions} & Liste filtrée, paginée et triée ; renvoie les données et les métadonnées. \\
GET & \texttt{/api/transactions/statistiques} & Agrégats par type et devise, avec les filtres de recherche. \\
GET & \texttt{/api/transactions/\{reference\}} & Détail d'une transaction. \\
POST & \texttt{/api/transactions} & Création d'une transaction de source \texttt{API} ; réponse 201 et en-tête \texttt{Location}. \\
PATCH & \texttt{/api/transactions/\{reference\}} & Modification partielle, après revalidation de l'objet complet. \\
DELETE & \texttt{/api/transactions/\{reference\}} & Suppression ; réponse 204 en cas de succès. \\
GET & \texttt{/api/rejets} et routes associées & Consultation des rejets, des statistiques par motif et du dernier chargement. \\
\bottomrule
\end{tabularx}
\caption{Principales routes de l'API}
\label{tab:routes-api}
\end{table}

\section{Filtres, pagination et tri}
La liste accepte des filtres par type, statut, canal, devise, catégorie, source, comptes, période et
montant. Les filtres d'énumération peuvent recevoir plusieurs valeurs séparées par une virgule.
Le paramètre \texttt{compte} recherche dans le compte source ou destination. La pagination utilise
\texttt{page} et \texttt{limite} (20 par défaut, 100 au maximum). Le tri est limité à la date, au
montant ou à la référence.

Un tri secondaire par référence départage les lignes ayant la même valeur de tri. Les pages restent
ainsi stables entre deux requêtes. Les paramètres inconnus ou invalides renvoient 400 au lieu d'être
ignorés silencieusement : une faute de frappe ne peut donc pas élargir accidentellement la recherche.

\section{Validation et réponses d'erreur}
Les champs sont normalisés (espaces retirés, codes en majuscules), puis contrôlés selon le contrat
commun décrit au chapitre~\ref{chap:donnees}. Les règles portent notamment sur les formats de
référence et de compte, les valeurs autorisées, les dates réelles, les montants positifs et les
cohérences entre champs. La catégorie de montant est calculée côté serveur. Les corps de requête
inconnus sont refusés.

Pour un \texttt{PATCH}, l'API fusionne les champs fournis avec la transaction existante, puis valide
l'ensemble. Un virement ne peut ainsi pas perdre son compte destination tout en restant accepté.
Les erreurs de validation sont regroupées par champ dans une réponse 400. Les principaux statuts
sont 400 (entrée invalide), 404 (ressource absente), 409 (référence déjà utilisée), 503 (base
indisponible) et 500 (erreur inattendue). Le format commun contient un code stable, un message et,
si nécessaire, le détail des champs fautifs.

\section{Sécurité des requêtes SQL}
Les valeurs reçues du client sont toujours transmises à PostgreSQL par paramètres \texttt{\$1},
\texttt{\$2}, etc. Elles ne sont pas concaténées au texte SQL. Comme les noms de colonnes ne
peuvent pas être paramétrés, les champs de tri sont choisis dans une liste blanche. Cette approche
réduit le risque d'injection SQL; elle ne remplace pas l'authentification ni les contrôles d'accès.

\section{Vérification}
Les tests unitaires couvrent la validation et la génération des requêtes paramétrées. Les tests
HTTP utilisent Supertest et un dépôt en mémoire pour contrôler les codes, les filtres, la
pagination, les écritures et les erreurs sans dépendre de PostgreSQL. Des tests d'intégration
complètent cette couverture sur une vraie base, et la collection Postman vérifie les parcours de
démonstration. Les résultats effectivement exécutés sont consignés dans le rapport de tests du
projet.